% Part: second-order-logic
% Chapter: syntax-and-semantics
% Section: expressive-power

\documentclass[../../../include/open-logic-section]{subfiles}

\begin{document}

\olfileid{sol}{syn}{exp}
\olsection{القوة التعبيرية}

\begin{explain}
تتسع قدرة اللغة على التعبير اتساعًا عظيمًا إذا جاز التكميم على متغيرات
الرتبة الثانية. فبالتكميم الوجودي فيها نقرر وجود دوال أو علاقات ذات
خصائص معينة، ولا يتأتى ذلك في منطق الرتبة الأولى إلا أن نخص الغرض
برمز غير منطقي، أي !!a{function} أو~!!{predicate}.
وبالتكميم الكلي من الرتبة الثانية نحكم على جميع المجموعات الجزئية من
!!{domain}، أو العلاقات عليه، أو الدوال منه إلى !!{domain}،
بأن لها خاصية بعينها. أما في الرتبة الأولى فليس لنا هذا العموم؛
وإنما نحكم على ما أُسند من المجموعات الجزئية والعلاقات والدوال إلى
رموز اللغة غير المنطقية. ثم إن الخاصية نفسها، التي نثبت وجود ما
يتصف بها أو نعممها على الجميع، يجوز أن يدخل في بيانها تكميم من
الرتبة الثانية. فمن هنا أمكن تعريف علاقات لا تُعرّف في الرتبة الأولى،
والتعبير عن خصائص للمجال تعجز الرتبة الأولى عن التعبير عنها.
\end{explain}

\begin{defn}
لتكن~$\Struct{M}$ !!a{structure} للغة~$\Lang{L}$.
نقول إن العلاقة~$R \subseteq \Domain{M}^2$ \emph{قابلة للتعريف}
في~$\Lang{L}$ إذا وجدت~!!{formula}~$!A_R(x, y)$ ليس لها من
المتغيرات الحرة إلا~$x$ و$y$، ويكون~$R(a, b)$، أي
$\tuple{a,b} \in R$، إذا وفقط إذا كان
$\Sat{M}{!A_R(x, y)}[s]$ حيث~$s(x) = a$ و$s(y) = b$.
\end{defn}

\begin{ex}
علاقة الهوية~$\Id{\Domain{M}}$، أي
$\Setabs{\tuple{a,a}}{a \in \Domain{M}}$، تُعرّف في منطق
الرتبة الأولى بالصيغة~$\eq[x][y]$. وأما في الرتبة الثانية فيمكن
تعريفها \emph{من غير~$\eq$}. وبيان ذلك أن~$a$ و$b$ إذا كانا
!!{element} واحدًا من~$\Domain{M}$ لزم أن يكونا !!{element}s
في المجموعات الجزئية نفسها من~$\Domain{M}$؛ إذ لا تتعين المجموعات
إلا بـ!!{element}s فيها. وإذا كان~$a$ و$b$ مختلفين، لم يكونا
!!{element}s في المجموعات الجزئية نفسها: فإن~$a \in \{a\}$
ولكن~$b \notin \{a\}$ متى كان~$a \neq b$.
فكونهما !!{element}s في المجموعات الجزئية نفسها
من~$\Domain{M}$ علاقة بين~$a$ و$b$ تتحقق إذا وفقط إذا
$a = b$. ولما كان التكميم على جميع المجموعات الجزئية
من~$\Domain{M}$ جائزًا في الرتبة الثانية، أمكن التعبير عن هذه
العلاقة فيها. فتكون !!{formula} الآتية تعريفًا
لـ$\Id{\Domain{M}}$:
\[
\lforall[X][(X(x) \liff X(y))]
\]
\end{ex}

\begin{prob}
أثبت أن~$\lforall[X][(X(x) \lif X(y))]$، وفيها~$\lif$
لا~$\liff$، تعرّف هي أيضًا~$\Id{\Domain{M}}$.
\end{prob}

\begin{ex}
ليكن~$R$ !!{predicate} ثنائي الموضع؛ فتكون~$\Assign{R}{M}$
علاقة ثنائية الموضع على~$\Domain{M}$. وسنستعمل الرمز~$R$
للدلالة على !!{predicate}~$R$ وعلى العلاقة~$\Assign{R}{M}$
نفسها، مع ما قد ينشأ عن جمع المعنيين في رمز واحد من التباس.
\emph{الإغلاق المتعدي}~$R^*$ للعلاقة~$R$ علاقة بين~$a$ و$b$
تتحقق إذا وفقط إذا وجدت~$c_1$، و\dots، و$c_k$ تتوالى بينها
العلاقات~$R(a,c_1)$، و$R(c_1, c_2)$، و\dots، و$R(c_k,b)$.
وندخل في ذلك حالة~$k = 0$؛ فصدق~$R(a,b)$ يستلزم
صدق~$R^*(a,b)$. ومن ثم~$R \subseteq R^*$.
بل إن~$R^*$ أصغر علاقة متعدية تحتوي~$R$.
وللتعبير في الرتبة الثانية عن كون~$X$ متعدية وتحتوي~$R$ نضع:
\begin{multline*}
  !B_R(X) \ident \lforall[x][\lforall[y][(R(x,y) \lif X(x, y))]] \land {}\\
\lforall[x][\lforall[y][\lforall[z][((X(x,y) \land X(y,z)) \lif X(x,
      z))]]].
\end{multline*}
فالمقترن الأول يقرر~$R \subseteq X$، والثاني يقرر تعدي~$X$.

ولكي تكون~$X$ أصغر العلاقات التي تحقق الشرطين، يلزم أن تكون
محتواة في كل علاقة متعدية تحتوي~$R$. وبذلك يُعرّف الإغلاق المتعدي
لـ$R$ بالـ!!{formula} الآتية:
\[
R^*(X) \ident !B_R(X) \land \lforall[Y][(!B_R(Y) \lif
  \lforall[x][\lforall[y][(X(x, y) \lif Y(x,y))]])].
\]
وعندئذ~$\Sat{M}{R^*(X)}[s]$ إذا وفقط إذا كان~$s(X) = R^*$.
وأما الإغلاق المتعدي لـ$R$ فلا يمكن التعبير عنه في منطق الرتبة
الأولى.
\end{ex}

\end{document}
